\chapter{Referencial Teórico}\label{cap:XX}

Normalmente o Capítulo~\ref{cap:XX} é parte do referencial teórico/estudo da arte/histórico... Aqui você precisa dar conhecimento para o leitor entender a sua proposta, deixando claro qual o problema. Pode ter mais de um capítulo no caso.


Continue usando citações como esta de exemplo~\cite{fangxing2010,gungor2011}.

Para itens estrangeiros como \textit{feedback}  não esqueça de colocar em itálico. Afirmações como essa: são os mesmos de 40 anos atrás~\cite{gungor2011}, precisam de referencia. 

\section{para seção}\label{sec:XX} %usar o label ajuda muito 

Usar o pacote de acronimos ajuda e muito, pois na primeira vez que o item aparecer no texto ele vai colocar o significado, e depois só a sigla ex: O \ac{SEP} tem como objetivo...

Para enumerar:

\begin{enumerate}
\item Geração: 
na geração, as usinas são classificadas conforme os recursos que utilizam, podendo ser hidroelétricas, termoelétricas, eólicas, nucleares, etc. Nesta fase, certo tipo de energia é transformado em energia elétrica. %A geração de energia inclui as instalações para esta finalidade. 
No caso do uso de hidrelétricas, envolve o armazenamento de um fluido, normalmente água de rios, conversão de energia hidráulica do fluido em energia mecânica em uma turbina hidráulica e a conversão da energia mecânica em energia elétrica por um gerador elétrico. 
\item Transmissão e Distribuição:
condução da energia elétrica das usinas de geração até os consumidores. Após gerada, a tensão é aumentada para que se evite perdas, em seguida, a energia elétrica é transportada em altas tensões, por fim, é distribuída ao consumidor final em baixas tensões. 
\end{enumerate}

Para siglas no plural use o acp: Em resumo, os \acp{TP} e \acp{TC} traduzem. 

Outro exemplo:  \acp{IED}, são capazes não só...

O label nas seções ajuda muito, veja: Conforme descrito na Seção~\ref{sec:XX}, ou no Capítulo~\ref{cap:XX}.

A Figura~\ref{fig:2_juncaomodelodados} precisa sempre vir referenciada.

\begin{figure}[ht]
	\begin{center}
		\centering
		\includegraphics[width=0.5\linewidth]{figuras/images}
	\end{center}	
	\caption{-----}
	\label{fig:2_juncaomodelodados}
\end{figure}

Exemplo de Figura lado a lado na Figura~\ref{fig:2_principioscom}.

\begin{figure*}[ht]
\begin{center}
\mbox{
	\subfigure[\ac{TPAA}, baseado em mensagens \textit{unicast}.]{\epsfig{figure=./figuras/images, width=0.4\textwidth}
	\label{subfig:2_principioscoma}
					}
	\hspace{0.2cm}
	\subfigure[\ac{MCAA}, baseado em mensagens \textit{multicast}.]{\epsfig{figure=./figuras/images, width=0.4\textwidth}
	\label{subfig:2_principioscomb}
					}
	}
\end{center}
\caption{Princípios de comunicação TPAA e MCAA da norma IEC 61850~\cite{fangxing2010}.}
\label{fig:2_principioscom}
\end{figure*}

E um exemplo de tabela, conforme Tabela ~\ref{tab:2_tabmsg}.

\begin{table}[h!]	
	\begin{center}
		\caption{Tipos de mensagens suportadas pelo padrão IEC 61850~\cite{fangxing2010}.}	
		%\setlength{\extrarowheight}{1.5pt}
		\begin{tabular}{|c|p{3.5cm}|p{4cm} |c|c|} \hline
		\label{tab:2_tabmsg}
{\bf Tipo} & {\bf Descrição} & {\bf Exemplo} & {\bf Classes} & {\bf Mensagem e }	\\
		   &			  &				  &	& Requisitos Temporais \\ \hline
		\hline
        1A & Mensagens Rápidas 				& {\it Trips} 					& P2 e P3 & GOOSE (3ms)		\\ \cline{4-5}
           &  				 				& 								& P1 	  & GOOSE (10ms)	\\ \hline 
        1B & Mensagens Rápidas 			 	& Comandos, 					& P2 e P3 & GOOSE (20ms)  	\\ \cline{4-5}
           & (outras)					 	& Mensagens Simples 			& P1	  & GOOSE (100ms) 	\\ \hline
         2 & Velocidade Média 				& Valores de Medidas 			& -  & MMS (100ms)			\\ \hline
         3 & Velocidade Baixa 				& Parâmetros 					& -  & MMS (500ms)			\\ \hline
         4 & Rajada de Dados                & Saída de dados dos 			& P2 e P3 & SV (3ms) 		\\ \cline{4-5}
		       &							&  transformadores				& P1      & SV (10ms)		\\ \hline
         5 & Transferência de Arquivos		& Arquivos grandes 				& -  & MMS($\geq$1000ms)	\\ \hline
        6A & Sincronização de Tempo A 		& Sincronização (\textit{station bus}) 		& -  & TimeSync	\\  \hline
        6B & Sincronização de Tempo A 		& Sincronização (\textit{process bus}) 		& -  & TimeSync \\ \hline
        7 & Mensagem de Comando 			& Comandos da estação HMI 		& -  & MMS(500ms)			\\ \hline
		\end{tabular}
	\end{center}
\end{table}

